% !TeX root = ../教材题库_全部条目_转写工作稿.tex
% Source: XJ-B1-C06-HW-P254-Y1370
% 第6章 / 6.4 用样本估计总体 / 教材 p.254、255 / PDF 第261、262页
% OCR未校对：公式与文字均可能有误；未完成数学TeX转写；禁止直接发给学生
% 原题及相邻原解截图（核对用）：../题目原图/XJ-B1-C06-HW-P254-Y1370-01.webp, ../题目原图/XJ-B1-C06-HW-P254-Y1370-02.webp
% 本文件只在转写完成后作为题干使用；OCR 条目尚不是可靠题干。
\begin{bankitem}{习题6.4 第3题}{第6章 / 6.4 用样本估计总体 / 教材 p.254、255 / PDF 第261、262页}{XJ-B1-C06-HW-P254-Y1370}{OCR未校对：公式与文字均可能有误；未完成数学TeX转写；禁止直接发给学生}
\begin{ocrdraft}
[教材p.254 / PDF第261页]
\par
3．抽样统计甲、
\par
乙两名射击运动员的5组（每组射击10次）训练成绩（单位：
\par
结果如下：
\par
运动员
\par
第一组
\par
第二组
\par
第三组
\par
第四组
\par
第五组
\par
甲
\par
91
\par
90
\par
89
\par
93
\par
乙
\par
89
\par
90
\par
91
\par
88
\par
92
\par
[教材p.255 / PDF第262页]
\par
从成绩的稳定性角度看，应派哪名运动员参加本周末的比赛比较合适？
\end{ocrdraft}
\end{bankitem}
